%%
%% fall-lecture06.tex
%% 
%% Made by Alex Nelson <pqnelson@gmail.com>
%% Login   <alex@lisp>
%% 
%% Started on  2026-10-13T07:55:03-0700
%% Last update 2026-10-13T07:55:03-0700
%% 

\lecture[]{}

\begin{node}
We are in the middle of deriving the equations of motion in
differential form. Recall, we derived the equations of motion for mass
(a.k.a., the ``continuity equation''):
\begin{equation}
\partial_{t}\rho+\partial_{i}(\rho u_{i})=0,
\end{equation}
and the conservation of linear momentum
\begin{equation}
\partial_{t}(\rho u_{i})+\partial_{j}(\rho u_{i}u_{j})-\gamma_{i}-\partial_{j}
\end{equation}
\end{node}

\begin{definition}
We say a fluid is \define{Inviscid} if
$\sigma_{ij}=-p\delta_{i,j}$. Physically this corresponds to no viscosity.
\end{definition}

\begin{example}[Bernoulli's principle for potential flows]
Assume steady flow (\S\ref{defn:steady-flow}), and we assume inviscid
flow. We assume body forces are the gradient of some potential
\begin{equation}
\gamma_{i}=-\rho\partial_{i}V.
\end{equation}
We further assume incompressibility
\begin{equation}
\rho=\mbox{constant}.
\end{equation}
The continuity equation becomes
\begin{equation}
\partial_{i}u_{i}=0.
\end{equation}
What happens to the conservation of linear momentum? We see, by steady flow
assumption,
\begin{equation}
\partial_{t}(\rho u_{i})=0
\end{equation}
which simplifies the conservation of linear momentum to (also pulling
density out in front of derivatives by incompressibility):
\begin{equation}
\rho\partial_{j}(u_{i}u_{j})-\gamma_{i}-\partial_{j}(-p\delta_{i,j})=0,
\end{equation}
then thhis simplifies further to
\begin{equation}
\rho\partial_{j}(u_{i}u_{j})+\rho\partial_{i}V+\partial_{i}p=0..
\end{equation}
Now the potential flow assumption:\footnote{This is just to expedite
the calculations, I do not believe it is physically necessary for
Bernoulli's principle.} there is some ``potential''
[scalar function] $\varphi$ such that
\begin{equation}
u_{i}=\partial_{i}\varphi.
\end{equation}
(Potential flow is equivalent to the velocity is irrotational.) We
plug this into the continuity equation to obtain
\begin{equation}
0=\partial_{i}u_{i}=\partial_{i}(\partial_{i}\varphi),
\end{equation}
which is just Laplace's equation
\begin{equation}
0=\nabla^{2}\varphi.
\end{equation}
The conservation of linear momentum equation becomes
\begin{equation}
\rho\partial_{j}\bigl((\partial_{i}\varphi)(\partial_{j}\varphi)\bigr)+\rho\partial_{i}V+\partial_{i}p=0.
\end{equation}
Then we use the product rule (of differentiation) and $\partial_{j}^{2}\varphi=0$
\begin{equation}
\rho(\partial_{j}\partial_{i}\varphi)(\partial_{j}\varphi)+\rho\partial_{i}V+\partial_{i}p=0.
\end{equation}
We assume
$\partial_{i}\partial_{j}\varphi=\partial_{j}\partial_{i}\varphi$ then
use the product rule again to find
\begin{equation}
\frac{1}{2}\rho\partial_{i}\bigr((\partial_{j}\varphi)(\partial_{j}\varphi)\bigr)+\rho\partial_{i}V+\partial_{i}p=0.
\end{equation}
Since $\rho$ is a constant, we can rewrite this as
\begin{equation}
\partial_{i}\left[\frac{\rho}{2}(\partial_{j}\varphi)(\partial_{j}\varphi)+\rho V+p\right]=0.
\end{equation}
Now, plugging $u_{j}=\partial_{j}\varphi$ back into this equation,
\begin{equation}
\partial_{i}\left[\frac{\rho}{2}u_{j}u_{j}+\rho V+p\right]=0.
\end{equation}
Then by the fundamental theorem of calculus,
\begin{equation}
\frac{\rho}{2}u_{j}u_{j}+\rho V+p=\mbox{constant}.
\end{equation}
This is precisely Bernoulli's principle.
\end{example}

\begin{definition}
We say a fluid is \define{Potential Flow} if there exists a scalar
function $\varphi$ such that the velocity vector field $u_{i}$ may be
written as $u_{i}=\partial_{i}\varphi$.
\end{definition}

\begin{node}[Conservation of energy]
Recall, we have
\begin{equation}
\varepsilon=e+\frac{1}{2}u_{i}u_{i},
\end{equation}
and the conservation of energy looks like
\begin{equation*}
\partial_{t}\int_{V_{c}}\rho\varepsilon\,\D V
+\oint_{\boundary V_{c}}\rho\varepsilon(\vec{u}\cdot\UnitNormalVec{n})\,\D A
=\int_{V_{c}}\vec{\gamma}\cdot\vec{u}\,\D V
+\oint_{\boundary V_{c}}\vec{f}\cdot\vec{u}\,\D A
-\oint_{\boundary V_{c}}\vec{q}\cdot\UnitNormalVec{n}\,\D A
+\int_{V_{c}}\rho Q\,\D V.
\end{equation*}
Then, using $f_{i}=\sigma_{ji}\widehat{n}_{j}$, using similar
reasoning as the other equations of motion,
\begin{equation}
\partial_{t}(\rho\varepsilon)+\partial_{j}(\rho u_{j}\varepsilon)
=\gamma_{j}u_{j}+\partial_{j}(u_{i}\sigma_{ji})-\partial_{j}q_{j}+\rho Q.
\end{equation}
This is 1 equation in 14 unknowns.
\end{node}

\subsection{Eulerian versus Lagrangian description of fluids}

\begin{node}
The Eulerian specification of a fluid flow has a fixed reference frame
(system of coordinates) and the fluid just ``flows along''. That is to
say, we look at the fluid flowing through the point $\vec{x}$.
In a fixed coordinate system, the observables look like $f(\vec{x},t)$.
\end{node}

\begin{node}
In the Lagrangian specification of a fluid flow, we follow a ``fluid parcel''
along its trajectory. (What is a ``fluid parcel''? Intuitively, it is
an ``infinitesimal volume'' of fluid which we can treat as a point
located at its center of mass.)
How do we specify a fluid parcel? We can use its initial position
$\vec{x}_{0}$ at time $t=t_{0}$ to specify the fluid parcel. The
trajectory is described using the function $\vec{x}(\vec{a},t)$ where
\begin{equation}
\vec{x}_{0}=\vec{x}(\vec{x}_{0},t_{0})
\end{equation}
is the initial position.

Note that we do not need to use the initial position $\vec{x}_{0}$ at
time $t_{0}$ as the first argument in the Lagrangian description of
fluid flow $\vec{x}(-,t)$. We could use \emph{any} parameter which
uniquely identifies each fluid parcel. For details about this, see
Lamb~\cite[\S16]{lamb1930hydrodynamics},
Batchelor~\cite[\S2.1]{batchelor1973introduction}.
\end{node}

\begin{node}
We can relate these two specifications using the velocity vector field
$\vec{u}(\vec{x},t)$ as
\begin{equation}
\vec{u}(\vec{x}(\vec{x}_{0},t),t)=\partial_{t}\vec{x}(\vec{x}_{0},t).
\end{equation}
This means any observable $f=f(\vec{x}(t),t)$ satisfies (by the chain rule):
\begin{subequations}
  \begin{align}
\frac{\D}{\D t}f 
&=\left.\partial_{t}f\right|_{\vec{x}~\text{fixed}}+\frac{\partial f}{\partial x_{j}}\frac{\partial x_{j}(t)}{\partial t}\\
&=\partial_{t}f|_{\vec{x}}+(\partial_{j}f)u_{j}\\
&=:\frac{D f}{D t}.
  \end{align}
\end{subequations}
For this reason, it is convenient to introduce the following.
\end{node}

\begin{definition}
Suppose a fluid follows the velocity vector field $\vec{u}(\vec{x},t)$
in the Eulerian specification. Then we define the \define{Material Derivative}
to be the differential operator
\begin{equation}
\frac{D}{D t}=\frac{\partial}{\partial t}+\sum_{j}u_{j}\frac{\partial}{\partial x_{j}}.
\end{equation}
In the Lagrangian specification of a fluid flow, this becomes ``just''
the total time derivative $\D/\D t$.
\end{definition}

\begin{xca}
Show the material derivative obeys the usual identities of the
derivative (the product rule and it is $\RR$-linear). Does it satisfy
the chain rule?
\end{xca}

\begin{example}
We see that
\begin{equation*}
\frac{D u_{i}}{D t}=\partial_{t}u_{i}+u_{j}\partial_{j}u_{i}.
\end{equation*}
\end{example}

\begin{example}
We see that
\begin{equation*}
\frac{D\rho}{D t}=\partial_{t}\rho + u_{j}\partial_{j}\rho.
\end{equation*}
\end{example}

\begin{node}[Continuity equation revisited]
We had the Continuity equation:
\begin{equation}
\partial_{t}\rho+\partial_{i}(u_{i}\rho)=0.
\end{equation}
We see that 
\begin{equation}
\partial_{t}\rho+\partial_{i}(u_{i}\rho)=\frac{D}{D t}\rho+\rho\partial_{i}u_{i}=0,
\end{equation}
or equivalently
\begin{equation}
\frac{1}{\rho}\frac{D\rho}{D t}=-\partial_{i}u_{i}.
\end{equation}
\end{node}

\begin{definition}
We define the \define{Dilitation} of a fluid flow to be the quantity
$\partial_{i}u_{i}$. 
\end{definition}

\begin{remark}
It is not immediately obvious, but it is equal to
the material derivative of the determinant of the Jacobian matrix of the Lagrangian specification
of the fluid flow:
\begin{equation*}
\frac{D}{D t}\det J(\vec{a},t)=(\vec{\nabla}\cdot\vec{u})\det\frac{\partial\vec{x}(\vec{a},t)}{\partial\vec{a}}.
\end{equation*}
See Childress~\cite[\S1.2]{childress2009introduction},
Lamb~\cite[\S6]{lamb1930hydrodynamics}, Batchelor~\cite[Eq~(2.2.3)~\textit{et~seq.}]{batchelor1973introduction}.
\end{remark}

\begin{node}
Let us consider some ``quantity'' $f(\vec{x},t)$. Then we see
\begin{subequations}
  \begin{align}
\partial_{t}(\rho f)+\partial_{j}\bigl((\rho u_{j})f\bigr)
&=\rho\partial_{t}f+f\partial_{t}\rho+\rho u_{j}\partial_{j}f+f\partial_{j}(\rho u_{j})\\
&=f[\partial_{t}\rho+\partial_{j}(\rho u_{j})]+\rho\frac{D f}{D t}\\
&=f[0] + \rho\frac{D f}{D t}\quad\mbox{by continuity equation}\\
&=\rho\frac{D f}{D t}.
  \end{align}
\end{subequations}
We see that the conservation of linear momentum has $f=u_{i}$, 
the conservation of energy had $f=\varepsilon$.
\end{node}

\begin{node}
The conservation of linear momentum equation is
\begin{equation*}
\rho\frac{D u_{i}}{D t}=\gamma_{i}+\partial_{j}\sigma_{ij}.
\end{equation*}
\end{node}

\begin{node}
The conservation of energy is then
\begin{equation*}
\rho\frac{D\varepsilon}{D t}=\gamma_{i}u_{i}+\partial_{j}(u_{i}\sigma_{ij})-\partial_{i}q_{i}+\rho Q.
\end{equation*}
\end{node}

\begin{remark}
References on the Eulerian and Lagrangian specifications which are also useful include Childress~\cite[Ch~1]{childress2009introduction}.
\end{remark}